\section{Sum Properties}
\label{sec:sum}

The recursive \texttt{sum} matches the mathematical summation, and
addition commutes over concatenation.

\begin{itemize}
  \item \href{https://github.com/thiagomata/prime-numbers/blob/master/src/main/scala/v1/chapter3/list/ListUtils.scala}{Sum matches summation}:
    $\operatorname{sum}(L) = L[0] + \cdots + L[|L|-1]$
  \item \href{https://github.com/thiagomata/prime-numbers/blob/master/src/main/scala/v1/chapter3/list/ListUtils.scala}{Left append preserves sum}:
    $\operatorname{sum}(x :: L) = x + \operatorname{sum}(L)$
  \item \href{https://github.com/thiagomata/prime-numbers/blob/master/src/main/scala/v1/chapter3/list/ListUtils.scala}{Sum over concatenation}:
    $\operatorname{sum}(A \mathbin{\texttt{++}} B) = \operatorname{sum}(A) + \operatorname{sum}(B)$
  \item \href{https://github.com/thiagomata/prime-numbers/blob/master/src/main/scala/v1/chapter3/list/ListUtils.scala}{Commutativity of sum}:
    $\operatorname{sum}(A \mathbin{\texttt{++}} B) = \operatorname{sum}(B \mathbin{\texttt{++}} A)$
\end{itemize}

\subsection{Sum matches Summation}

We can prove that the recursive \texttt{sum} function over a list $L$
matches the mathematical definition of the summation
$\sum_{i=0}^{n-1} x_i$, where $L = [x_0, x_1, \dots, x_{n-1}]$, $|L| = n$.

\textbf{Base case}: $|L| = 0$

\begin{equation*}
\begin{aligned}
\operatorname{sum}(L) &= 0 & \text{[by definition of sum]} \\
\sum L &= 0 & \text{[summation over empty list]} \\
\implies \operatorname{sum}(L) &= \sum L \in \mathbb{L}
\end{aligned}
\end{equation*}

\begin{equation*}
\therefore
\end{equation*}

\begin{equation*}
\forall\, L \in \mathbb{L}
\end{equation*}

\begin{equation*}
|L| = 0 \implies \operatorname{sum}(L) = \sum L
\end{equation*}

\textbf{Inductive step}: $|L| > 0$

Let $P \in \mathbb{L}$, with $P = [x_1, x_2, \dots, x_{n-1}] \in \mathbb{L}$,
and assume:

\begin{equation*}
\begin{aligned}
\operatorname{sum}(P) & = \sum_{i=1}^{n-1} x_i && \text{[by Inductive Hypothesis]} \\
L = x_0 :: P & = [x_0, x_1, \dots, x_{n-1}] && \text{[by Definition of Cons]} \\
\end{aligned}
\end{equation*}

We can ensure termination, since:

\begin{equation*}
\begin{aligned}
|L| &= |P| + 1 && \text{[by Size Definition]} \\
|P| &< |L| && \text{[Size Decreases Ensures Termination]} \\
\end{aligned}
\end{equation*}

Let's calculate the sum of $L$:

\begin{equation*}
\begin{aligned}
\operatorname{sum}(L) &= \operatorname{head}(L) + \operatorname{sum}(\operatorname{tail}(L)) && \text{[by definition of the recursive function sum]} \\
&= x_0 + \operatorname{sum}(P) && \text{[by definition of head and P]} \\
&= x_0 + \sum_{i=1}^{n-1} x_i && \text{[by Inductive Hypothesis]} \\
&= \sum_{i=0}^{n-1} x_i = \sum L
\end{aligned}
\end{equation*}

\begin{equation*}
\therefore
\end{equation*}

\begin{equation*}
\forall\, L \in \mathbb{L}
\end{equation*}

\begin{equation*}
|L| > 0 \implies \operatorname{sum}(L) = \sum L
\end{equation*}

Hence, by induction on the size of $L$:

\begin{equation*}
\forall\, L \in \mathbb{L}
\end{equation*}

\begin{equation*}
\operatorname{sum}(L) = \sum L = \sum_{i=0}^{n-1} x_i \in \mathbb{S} \quad \blacksquare \quad \text{[Q.E.D.]}
\end{equation*}

This property is verified in
\href{https://github.com/thiagomata/prime-numbers/blob/master/src/main/scala/v1/chapter3/list/ListUtils.scala}{\texttt{ListUtils::sum}}.
The implementation and verification code are in Appendix~A.7.

\subsection{Left Append Preserves Sum}

The sum of a list with an element prepended equals the element plus the
sum of the original list.

\begin{equation*}
\forall\, x \in \mathbb{S}
\end{equation*}

\begin{equation*}
\operatorname{sum}(x :: L) = x + \operatorname{sum}(L)
\end{equation*}

Proof:

\begin{equation*}
\begin{aligned}
A & = x :: L && \text{[Cons]} \\
\operatorname{sum}(A) & = \operatorname{head}(A) + \operatorname{sum}(\operatorname{tail}(A)) && \text{[By recursive definition of sum]} \\
& = x + \operatorname{sum}(L) && \text{[By recursive definition of head and tail]} \\
\end{aligned}
\end{equation*}

\begin{equation*}
\therefore
\end{equation*}

\begin{equation*}
\begin{aligned}
\operatorname{sum}(x :: L) & = x + \operatorname{sum}(L) \quad \blacksquare && \text{[Q.E.D.]} \\
\end{aligned}
\end{equation*}

{\raggedright
This property is verified in
\href{https://github.com/thiagomata/prime-numbers/blob/master/src/main/scala/v1/chapter3/list/ListUtils.scala}{\texttt{ListUtils::\allowbreak listSumAddValue}}.
The full Scala verification code is in Appendix~A.8.
\par}

\subsection{Sum over Concatenation}

The sum of two concatenated lists equals the sum of each list added
together.

\begin{equation*}
\operatorname{sum}(A \mathbin{\texttt{++}} B) = \operatorname{sum}(A) + \operatorname{sum}(B)
\end{equation*}

\textbf{If list A is empty}:

\begin{equation*}
\begin{aligned}
A \mathbin{\texttt{++}} B & = L_e \mathbin{\texttt{++}} B & \text{[A is empty list]} \\
& = B & \text{[By definition of concatenation]} \\
\operatorname{sum}(A) & = 0 & \text{[By definition of sum]} \\
\operatorname{sum}(A \mathbin{\texttt{++}} B) & = \operatorname{sum}(B) & \text{[Since A} \mathbin{\texttt{++}} \text{B equals B]} \\
& = 0 + \operatorname{sum}(B) \\
& = \operatorname{sum}(A) + \operatorname{sum}(B) & \text{[Since sum(A) is zero]} \\
\end{aligned}
\end{equation*}

\textbf{If list A is non-empty}:

\begin{equation*}
\begin{aligned}
C & = \operatorname{tail}(A) \mathbin{\texttt{++}} B \\
\operatorname{sum}(A) & = \operatorname{head}(A) + \operatorname{sum}(\operatorname{tail}(A)) & \text{[By definition of sum]} \\
\operatorname{sum}(C) & = \operatorname{sum}(\operatorname{tail}(A)) + \operatorname{sum}(B) & \text{[Inductive Step]} \\
A \mathbin{\texttt{++}} B & = \operatorname{head}(A) :: (\operatorname{tail}(A) \mathbin{\texttt{++}} B) & \text{[By definition of head and tail]} \\
\operatorname{sum}(A \mathbin{\texttt{++}} B) & = \operatorname{head}(A) + \operatorname{sum}(\operatorname{tail}(A) \mathbin{\texttt{++}} B) & \text{[By definition of sum]} \\
& = \operatorname{head}(A) + \operatorname{sum}(\operatorname{tail}(A)) + \operatorname{sum}(B) & \text{[By definition of C]} \\
& = \operatorname{sum}(A) + \operatorname{sum}(B) & \text{[Substituting]} \\
\end{aligned}
\end{equation*}

\begin{equation*}
\therefore
\end{equation*}

\begin{equation*}
\begin{aligned}
\operatorname{sum}(A \mathbin{\texttt{++}} B) = \operatorname{sum}(A) + \operatorname{sum}(B) & \quad \blacksquare \qquad \text{[Q.E.D.]} \\
\end{aligned}
\end{equation*}

{\raggedright
This property is verified in
\href{https://github.com/thiagomata/prime-numbers/blob/master/src/main/scala/v1/chapter3/list/ListUtils.scala}{\texttt{ListUtils::\allowbreak listCombine}}.
The full Scala verification code is in Appendix~A.9.
\par}

\subsection{Commutativity of Sum over Concatenation}

The order of concatenation does not affect the total sum.

\begin{equation*}
\operatorname{sum}(A \mathbin{\texttt{++}} B) = \operatorname{sum}(B \mathbin{\texttt{++}} A)
\end{equation*}

Since:

\begin{equation*}
\begin{aligned}
\operatorname{sum}(A \mathbin{\texttt{++}} B) & = \operatorname{sum}(A) + \operatorname{sum}(B) & \text{[Sum over Concatenation]} \\
\operatorname{sum}(B \mathbin{\texttt{++}} A) & = \operatorname{sum}(B) + \operatorname{sum}(A) & \text{[Sum over Concatenation]} \\
\operatorname{sum}(B) + \operatorname{sum}(A) & = \operatorname{sum}(A) + \operatorname{sum}(B) & \text{[Distributive]} \\
\operatorname{sum}(B \mathbin{\texttt{++}} A) & = \operatorname{sum}(A \mathbin{\texttt{++}} B) \quad \blacksquare & \text{[Q.E.D]} \\
\end{aligned}
\end{equation*}

{\raggedright
This property is verified in
\href{https://github.com/thiagomata/prime-numbers/blob/master/src/main/scala/v1/chapter3/list/ListUtils.scala}{\texttt{ListUtils::\allowbreak listSwap}}.
The full Scala verification code is in Appendix~A.10.
\par}

\subsection{Sum Positivity}

If every element of a non-empty list is greater than zero, the sum of the
list is greater than zero.

\begin{equation*}
\begin{aligned}
(\forall x \in L,\, x > 0) \wedge L \neq L_e \implies \operatorname{sum}(L) > 0
\end{aligned}
\end{equation*}

The base case is the singleton list, where the sum is just the one
positive element. The inductive step adds a positive head to a tail whose
sum is already known to be positive by the inductive hypothesis.

{\raggedright
This property is verified in
\href{https://github.com/thiagomata/prime-numbers/blob/master/src/main/scala/v1/chapter3/list/properties/ListUtilsProperties.scala}{\texttt{ListUtilsProperties::\allowbreak assertSumPositive}}.
\par}
